%% ma: L11-B09-0009
%% lop: 11
%% chuong: 3
%% bai: 9
%% dang: 11.9.4
%% loai: TLN
%% muc: TH
%% dap_an: "94,2"
%% nguon: "de-on-tap-giua-ki-1-toan-11-nam-2026-2027-truong-thpt-ngo-quyen-da-nang.pdf (M09, Đề ôn tập giữa kì 1 Toán 11 THPT Ngô Quyền, Đà Nẵng), Mã đề 0111, Phần III Câu 4"
%% kien_thuc: "Bài 9 (các số đặc trưng đo xu thế trung tâm)"
%% trang_thai: chinh-thuc
\begin{ex}
Khảo sát thời gian sử dụng Internet mỗi ngày (đơn vị: phút) của $50$ học sinh được ghi lại ở bảng số liệu ghép nhóm sau:
\begin{center}\begin{tabular}{|l|c|c|c|c|c|}\hline
Thời gian (phút) & $[30;60)$ & $[60;90)$ & $[90;120)$ & $[120;150)$ & $[150;180)$\\\hline
Số học sinh & $8$ & $15$ & $17$ & $7$ & $3$\\\hline
\end{tabular}\end{center}
Tính thời gian sử dụng Internet trung bình mỗi ngày của $50$ học sinh trên (đơn vị: phút).
\shortans{94,2}
\loigiai{Giá trị đại diện: $45;75;105;135;165$. \[\bar x=\dfrac{45\cdot8+75\cdot15+105\cdot17+135\cdot7+165\cdot3}{50}=\dfrac{4710}{50}=94{,}2\text{ (phút).}\]}
\end{ex}
